%=============================================================================
% ABSTRACT
%=============================================================================
\begin{abstract}
Spatio-temporal graph neural networks (ST-GNNs) are widely applied to
district-level epidemic forecasting, yet their advantage over naive temporal
baselines is rarely verified under a leak-free protocol. We reproduce five
published ST-GNNs on Sri Lanka's national dengue record (25 districts, 459 weeks)
to within 7.3\% across twenty comparisons, and find that the reported gains are
measured with training data included; naive persistence beats every model in the
column the source paper reports. Re-evaluating under rolling-origin
cross-validation, we identify an administrative reporting backlog at week~395 that
occupies one fold's test split and carries $\sim$90\% of its squared error,
inverting which fold appears hardest. Scoring every arm with and without those
windows, the persistence floor moves from RMSE 44.80 to 29.52, and the
best-looking architecture's 3.8-point lead disappears. Against the artifact-free
floor no architecture and none of four evidence-led increments differ from their
own baseline at $p<0.05$. We then test the mechanistic prior the task invites: a
renewal equation whose generation interval is derived from a validated SEIR--SEI
model. It reconstructs the record almost exactly given oracle $R_t$ (RMSE 13.44)
and fails as a forecaster, because $R_t$ is only 26\% predictable out of sample,
so routing a forecast through it re-injects a 1.90$\times$ multiplicative error
where persistence carries none. The ceiling is informational rather than
architectural: the models' predicted growth carries no information about realised
growth (regression slope 0.002), and under squared error that is optimal. The
mechanistic quantity is nonetheless the single largest contributor to a different
and tractable target---ranking which districts will exceed an alert threshold
three weeks ahead (AUC $0.807 \rightarrow 0.826$).
\end{abstract}
